\documentclass[10pt,a4paper]{amsart}
\usepackage[margin=19mm]{geometry}
\usepackage{amsmath,amssymb,mathtools}
\usepackage[scr=rsfs,cal=euler]{mathalfa}
\usepackage{xfrac}
\usepackage[unicode,hidelinks,bookmarks=false]{hyperref}
\newcommand{\ie}{i.e.\ }
\newcommand{\cf}{cf.\ }
\newcommand{\resp}{respectively\ }
\newcommand{\Subsection}{\subsection}
\providecommand{\nobreakdash}{\nobreak\hbox{-}}
\setlength{\emergencystretch}{2em}
\multlinegap0pt
\usepackage{amssymb}
\usepackage{xcolor}
\usepackage{algorithm}\usepackage{algpseudocode}
\usepackage[shortlabels]{enumitem}
\usepackage{float}
\newtheorem{theorem}{Theorem}
\title{Computational Oncology of Chemotaxis-Driven Tumour--Immune Spatial Patterning and Stability}
\author{Zonghao Liu, Jiguang Yu, Lei Su, Louis Shuo Wang, Yang Du, Jingfeng Liu}
\date{Source: arXiv:2607.03813v2 (CC BY 4.0)}
\begin{document}
\maketitle

\noindent\emph{Source and licence.} Computational Oncology of Chemotaxis-Driven Tumour--Immune Spatial Patterning and Stability. Version arXiv:2607.03813v2 (CC BY 4.0), distributed by arXiv under the Creative Commons Attribution 4.0 International licence, http://creativecommons.org/licenses/by/4.0/, as recorded in the arXiv licence metadata for this exact version and shown on the abstract page https://arxiv.org/abs/2607.03813v2. Full licence notice: \texttt{licenses/arxiv-2607.03813-CC-BY-4.0.txt}.\par\smallskip

\noindent\textit{Editorial scope.} Counted unit: the article's theorem thm:tumour-free-threshold (source0.tex lines 1258--1266). The article's complete original proof of it is delivered with it (source0.tex lines 1268--1321, one \texttt{proof} environment of 54 lines). No mathematics is written, repaired or re-derived: the statement and the proof are the article's own lines, in the article's own notation, and no retained line is substituted or re-flowed.  Retained uncounted as the source-local context the proof needs: lines 550--557.  Nothing was omitted from any retained span, so the package carries no omission record.  No retained line contains a citation, so the wrapper carries no bibliography and no bibliography entry.  No retained line contains a figure environment, an image-inclusion command or drawing code, so no image is delivered and the package declares no asset dependency.  Editorial formatting only, in addition: an AMS \texttt{amsart} preamble that restates the article's own declarations and macros used by the retained lines, namely the environments \texttt{theorem} on separate counters, together with the standard packages those lines need.  The retained environments print in this document as Theorem 1 in order of appearance, whereas the article prints the same items with the numbering of its own full document.  The complete native source of the article as distributed in the arXiv e-print for this exact version is bundled unchanged as \texttt{sources/204424/source0.tex}.
\begin{equation}\label{eq:nondim-model}
\begin{aligned}
u_t &= d_1\Delta u+u(1-u-v),\\
v_t &= d_2\Delta v-\xi\nabla\cdot(v\nabla w)
      +\sigma_0+\sigma_1\frac{w}{1+w}-\delta v-\beta uv,\\
w_t &= d_3\Delta w+\alpha u+\gamma uv-\ell w.
\end{aligned}
\end{equation}

\begin{theorem}[Tumour-free threshold]
\label{thm:tumour-free-threshold}
The tumour-free equilibrium $E_0$
is linearly stable if $\sigma_0>\delta$,
linearly unstable if $\sigma_0<\delta$,
and non-hyperbolic at $\sigma_0=\delta$.
At this equilibrium, chemotaxis does not affect the linear tumour-invasion
threshold.
\end{theorem}

\begin{proof}
Let
\[
u=0+\widetilde u,\qquad
v=\frac{\sigma_0}{\delta}+\widetilde v,\qquad
w=0+\widetilde w.
\]
The linearization of \eqref{eq:nondim-model} at $E_0$ is
\begin{equation}\label{eq:tumour-free-linearization}
\begin{cases}
\widetilde u_t=
d_1\Delta \widetilde u+
\left(1-\frac{\sigma_0}{\delta}\right)\widetilde u,\\[2pt]
\widetilde v_t=
d_2\Delta \widetilde v-\delta\widetilde v
+\sigma_1\widetilde w
-\xi\frac{\sigma_0}{\delta}\Delta\widetilde w
-\beta\frac{\sigma_0}{\delta}\widetilde u,\\[2pt]
\widetilde w_t=
d_3\Delta \widetilde w+
\left(\alpha+\gamma\frac{\sigma_0}{\delta}\right)\widetilde u
-\ell\widetilde w .
\end{cases}
\end{equation}
Using the Neumann Laplacian eigenbasis
\[
-\Delta\phi_k=\mu_k\phi_k,\qquad
\partial_n\phi_k=0,
\qquad
0=\mu_0<\mu_1\le\mu_2\le\cdots,
\]
the tumour component has growth rate
$\lambda_k^{(u)}
=
1-\dfrac{\sigma_0}{\delta}-d_1\mu_k$.
The largest of these rates occurs at $k=0$ and equals
$\lambda_0^{(u)}=1-\dfrac{\sigma_0}{\delta}$.
Hence all tumour perturbations decay if $\sigma_0>\delta$, while the homogeneous
tumour mode grows if $\sigma_0<\delta$.

The remaining two components form a triangularly forced stable subsystem once
$\widetilde u$ is fixed. Indeed, the diagonal decay rates for
$\widetilde v$ and $\widetilde w$ contain $-\delta-d_2\mu_k$ and
$-\ell-d_3\mu_k$, which are strictly negative. Therefore, the sign of the
dominant eigenvalue is determined by
$1-\sigma_0/\delta$. This proves the stated stability threshold.

Finally, the chemotactic contribution appears only in the immune equation
through the linear term
$-\xi\dfrac{\sigma_0}{\delta}\Delta\widetilde w$.
It does not enter the tumour growth rate
$\lambda_k^{(u)}$. Thus chemotaxis does not affect the linear
tumour-invasion threshold at $E_0$.
\end{proof}

\end{document}
